\documentclass[acmsmall]{acmart}

\citestyle{acmauthoryear}

\title{Towards End-to-End Quality Assessment of LLM-Based Systems:
A Survey of Quality Dimensions, Evaluation Methods, and Integration}

% TODO: replace with the final author list and affiliations.
\author{Dener Silva}

\begin{abstract}
Large Language Models (LLMs) have evolved from standalone text generators into components of increasingly complex software systems that combine prompt engineering, retrieval-augmented generation, tools, memory, multi-agent coordination, and production infrastructure. This evolution has produced an equally fragmented evaluation landscape: different studies assess prompt robustness, retrieval quality, factuality, task success, judge agreement, latency, safety, or cost, often in isolation and with limited visibility into how quality propagates across the complete system lifecycle.

This survey organizes LLM quality assessment around two complementary questions: \emph{what} aspect of an LLM-based system is evaluated and \emph{how} that quality is measured. We distinguish quality dimensions associated with input and specification, context and retrieval, generated output, conversational interaction, agent processes and trajectories, and operational/system behavior. Separately, we organize evaluation mechanisms including deterministic and reference-based metrics, semantic evaluators, learned evaluators, human evaluation, LLM-as-a-Judge, multi-judge approaches, task- and environment-based evaluation, adversarial testing, perturbation-based methods, trace-aware evaluation, and statistical reliability analysis.

Beyond cataloguing individual metrics, we investigate the degree to which existing approaches integrate multiple stages of the LLM lifecycle. We distinguish isolated evaluation, cross-stage assessment, black-box end-to-end evaluation, and lifecycle-integrated evaluation. Evaluator reliability is treated as a meta-evaluation problem rather than as a normal lifecycle stage, allowing the use of an LLM judge to be distinguished from evidence that the judge itself is reliable. The resulting synthesis is intended to clarify how current evaluation practices relate to one another and to identify gaps toward more coherent, traceable, and end-to-end quality assessment of LLM-based systems.
\end{abstract}

\keywords{Large Language Models, LLM Evaluation, Quality Assessment,
LLM-as-a-Judge, Retrieval-Augmented Generation, AI Agents,
End-to-End Evaluation, Reliability}

\begin{document}

\maketitle

\section{Introduction}
\label{sec:introduction}

Large Language Models (LLMs) have evolved from standalone text generation models into components of increasingly complex software systems. Modern LLM-based applications may incorporate prompt templates, retrieval-augmented generation (RAG), external tools, memory mechanisms, multi-agent coordination, iterative reasoning, and interaction with dynamic environments. Consequently, the quality of an LLM-based system can no longer be adequately characterized only by measuring the linguistic quality or factual correctness of its final response.

A failure observed at the output of an LLM system may originate at several different stages. An incorrect response may result from an ambiguous or incomplete user specification, inadequate retrieved context, an error in reasoning or tool selection, a failure during a multi-step agent trajectory, or limitations of the underlying model itself. Similarly, a response that appears correct according to an automated evaluator may still be incorrectly assessed if the evaluation mechanism is biased, unstable, poorly calibrated, or insufficiently aligned with human judgment. In production environments, additional factors such as latency, cost, availability, throughput, resource consumption, privacy, and safety may further determine whether a technically capable system can be considered satisfactory in practice.

This complexity has resulted in a rapidly expanding evaluation ecosystem. Numerous benchmarks and metrics have been proposed to assess individual properties of LLMs and LLM-based systems, while evaluation frameworks increasingly combine automated metrics, human feedback, LLM-based judges, execution-based tests, traces, and operational telemetry. At the same time, different research communities often employ distinct terminology for closely related concepts, and apparently similar metrics may operationalize different notions of quality.

A particularly visible example is the growing use of LLM-as-a-Judge approaches. LLM judges have become a common \emph{method} for evaluating model outputs, agent trajectories, and other artifacts. However, using an LLM as an evaluator is conceptually different from evaluating the quality of that evaluator. Reliability, consistency, bias, calibration, agreement with human experts, and sensitivity to evaluation prompts constitute a separate meta-evaluation problem. Conflating these two roles can obscure an important source of uncertainty in automated evaluation.

A similar distinction arises in end-to-end evaluation. A system may be evaluated from input to final task outcome without examining its internal stages. Such black-box evaluation is useful for measuring overall task success, but it provides limited information about the origin of failures. In contrast, other approaches explicitly relate quality across multiple stages, such as retrieval quality and response quality, or agent trajectory quality and task completion. A smaller subset attempts to combine several lifecycle stages together with evaluator reliability or operational characteristics. These approaches represent qualitatively different levels of evaluation integration.

This survey investigates this fragmented landscape from a system-oriented perspective. Rather than organizing the literature around individual benchmarks or particular evaluation technologies, we consider quality assessment across the lifecycle of an LLM-based system. Our central premise is that a comprehensive understanding of LLM quality requires distinguishing three complementary aspects: the object being evaluated, the quality dimension being measured, and the mechanism used to perform that measurement.

Accordingly, this work makes the following contributions:

\begin{itemize}
    \item We propose a lifecycle-oriented taxonomy of quality assessment for LLM-based systems, covering input and specification, context and retrieval, output, conversational interaction, agent processes and trajectories, and operational/system quality.
    \item We distinguish the \emph{quality dimension} being evaluated from the \emph{evaluation method} used to assess it, organizing methods such as deterministic metrics, semantic metrics, human evaluation, LLM-as-a-Judge, multi-judge evaluation, task-based evaluation, adversarial testing, trace-aware evaluation, and statistical reliability analysis.
    \item We explicitly model evaluator quality as a meta-evaluation problem, separating the use of an evaluator from evidence regarding its reliability.
    \item We introduce a taxonomy of evaluation integration that distinguishes isolated evaluation, cross-stage evaluation, black-box end-to-end evaluation, and lifecycle-integrated evaluation.
    \item We analyze how existing approaches combine quality dimensions, evaluation methods, and operational considerations, with the goal of identifying weakly connected areas despite the large number of existing metrics and frameworks.
\end{itemize}

The remainder of this article is organized as follows. Section~\ref{sec:background} introduces the conceptual foundations and terminology used throughout the survey. Section~\ref{sec:research-questions} presents the research questions. Section~\ref{sec:methodology} describes the literature search, eligibility criteria, classification protocol, and human validation procedure. Section~\ref{sec:taxonomy} presents the proposed taxonomy. The subsequent evidence-driven sections will synthesize evaluation methods, integration patterns, tools, and research gaps after completion of the corpus audit.

\section{Background and Scope}
\label{sec:background}

\subsection{From LLM Evaluation to LLM-System Evaluation}

Traditional language model evaluation primarily focuses on the relationship between an input and a model output. This abstraction remains useful for many tasks, but it is increasingly insufficient for modern LLM applications. An LLM-based system may contain several interacting components before a final result is produced:

\begin{equation}
x \rightarrow S \rightarrow C \rightarrow P \rightarrow y \rightarrow E,
\end{equation}

where $x$ represents the original user input, $S$ represents its specification or prompt construction, $C$ represents retrieved or externally supplied context, $P$ denotes the execution process---including reasoning, tools, agents, or trajectories---$y$ represents the resulting output or task outcome, and $E$ denotes the evaluation process.

This representation is intentionally general. Not every LLM-based system contains all stages, and the boundaries between stages vary across architectures. Its purpose is to highlight that observed output quality may depend on several upstream components and that the evaluation mechanism itself introduces an additional layer that may require validation.

\subsection{Quality Dimensions and Evaluation Methods}

A central distinction in this survey is between \emph{what is evaluated} and \emph{how it is evaluated}. For example, the faithfulness of a RAG response is a quality property, whereas an LLM judge used to score that faithfulness is an evaluation method. Likewise, agent task completion is a quality outcome, while execution tests or state matching are evaluation mechanisms.

We therefore organize the literature along two orthogonal axes. The first axis represents the primary quality dimensions:

\begin{enumerate}
    \item \textbf{Input and Specification Quality};
    \item \textbf{Context and Retrieval Quality};
    \item \textbf{Output Quality};
    \item \textbf{Conversation and Interaction Quality};
    \item \textbf{Agent, Process, and Trajectory Quality};
    \item \textbf{Operational and System Quality}.
\end{enumerate}

Evaluator or judge quality is treated separately as a meta-evaluation dimension because it characterizes the reliability of the measurement mechanism rather than a normal stage of the LLM execution lifecycle.

The second axis represents evaluation methods. The literature includes deterministic and rule-based checks, lexical and reference-based metrics, semantic metrics, learned evaluators, human evaluation, LLM-as-a-Judge, multi-judge or ensemble approaches, task- and environment-based verification, adversarial and red-team testing, perturbation-based evaluation, trace-aware methods, and statistical reliability analysis.

\subsection{Quality Target, Quality Dimension, and Evaluation Method}

Three concepts are kept separate throughout the review. The \emph{evaluation object} denotes what is primarily being assessed, such as an LLM-based system, an evaluator/judge, or an external-domain object. The \emph{quality dimension} denotes the property or lifecycle stage whose quality is under investigation. The \emph{evaluation method} denotes the mechanism used to produce a score, judgment, ranking, or verification outcome.

This distinction is particularly important for domain-specific studies. The presence of an LLM in a healthcare, energy, IoT, or software-engineering pipeline does not automatically imply that LLM quality is being evaluated. Conversely, a domain-specific study may be directly relevant when it independently assesses the quality contribution of the LLM-based subsystem.

\section{Research Questions}
\label{sec:research-questions}

The survey is guided by six research questions designed to move from descriptive mapping to integration analysis and gap identification.

\begin{description}

\item[RQ1:] \textbf{Which quality dimensions of LLM-based systems are evaluated in the existing literature?}

RQ1 maps the literature onto the lifecycle-oriented taxonomy and examines how attention is distributed across input/specification, context/retrieval, output, conversation/interaction, agent/process/trajectory, and operational/system quality. It also records cross-cutting concerns such as robustness, safety, security, fairness, multimodality, and efficiency.

\item[RQ2:] \textbf{Which evaluation methods and protocols are used to assess these quality dimensions?}

RQ2 separates the construct being evaluated from the mechanism used to measure it. We examine deterministic and rule-based checks, lexical and reference-based metrics, semantic metrics, learned evaluators, human evaluation, LLM-as-a-Judge, multi-judge approaches, task- and environment-based verification, adversarial testing, perturbation-based methods, trace-aware evaluation, and statistical reliability analysis. We additionally record protocols such as absolute scoring, pairwise comparison, ranking, rubric scoring, binary classification, critique, preference judgment, and task-success verification.

\item[RQ3:] \textbf{To what extent does the literature connect quality across multiple stages of an LLM-based system?}

RQ3 investigates whether evaluation remains isolated at one component, explicitly relates multiple stages, measures a multi-stage workflow only through its final outcome, or provides a lifecycle-integrated assessment. This question directly addresses the difference between black-box end-to-end evaluation and stage-aware quality attribution.

\item[RQ4:] \textbf{How is the reliability of the evaluation mechanism itself assessed?}

RQ4 treats evaluator quality as a meta-evaluation problem. We examine whether studies assess consistency, human agreement, calibration, bias, robustness, prompt sensitivity, dependence among multiple evaluators, or other properties that determine whether an automated or human judgment can be trusted.

\item[RQ5:] \textbf{How are operational and deployment-oriented properties incorporated into LLM quality assessment?}

RQ5 examines whether semantic or task-level quality is considered together with latency, throughput, monetary cost, token and memory consumption, energy use, availability, observability, privacy, safety, and related system-level constraints. The goal is to identify whether operational quality is treated as part of overall system quality or as a separate engineering concern.

\item[RQ6:] \textbf{What evidence gaps remain toward reliable end-to-end quality assessment of LLM-based systems?}

RQ6 synthesizes the preceding results to identify dimensions that are well studied in isolation but weakly connected, evaluation methods whose reliability is insufficiently characterized, and lifecycle relationships for which empirical evidence remains scarce. This final question provides the basis for the research agenda proposed in the survey.

\end{description}

\section{Review Methodology}
\label{sec:methodology}

\subsection{Review Design}

This study follows a structured literature review process combining broad literature retrieval, automated support for initial screening and classification, targeted full-text evidence extraction, and independent human validation of the classification protocol. The objective is not only to identify evaluation metrics, but also to characterize the object of evaluation, quality dimensions, evaluation mechanisms, and the degree of integration between system stages.

The historical retrieval phase produced 2,662 records across 12 search blocks. Because those records precede the final publication-status audit, this number represents the \emph{retrieved historical corpus}, not the final set of peer-reviewed studies included in the academic synthesis. The historical corpus is frozen at May 25, 2026, and a temporal update is used to capture relevant publications appearing after that date.

\subsection{Eligibility Criteria}

A study is considered \textbf{Core} when it directly evaluates one or more quality dimensions of an LLM or LLM-based system, or explicitly investigates relationships between such dimensions. A study is considered \textbf{Adjacent} when an LLM is primarily used as a tool or component to evaluate an external-domain object, or when its methodology is potentially transferable to LLM quality assessment but LLM-system quality is not the principal object of investigation. Studies without substantive relevance to LLM-system quality assessment are excluded.

The presence of an external application domain does not by itself imply an Adjacent classification. A healthcare, energy, IoT, or software-engineering study may be Core if it directly measures the quality of the underlying LLM-based system. Conversely, a study in which an LLM is merely used to estimate or classify properties of an external system is classified as Adjacent unless the LLM contribution is separately evaluated.

\subsection{Peer-Review Eligibility}

Peer review is treated as an explicit eligibility criterion for the primary academic corpus. The existence of an arXiv version does not imply exclusion when a peer-reviewed version of record exists. Conversely, preprint-only works, technical reports, theses, white papers, and other non-peer-reviewed material are excluded from the primary academic counts, although relevant material may be retained in a separate supplementary or practice-oriented evidence track.

For secondary studies, we additionally record whether their underlying evidence base is restricted to peer-reviewed literature or mixes peer-reviewed studies with preprints and other grey literature. This distinction prevents the publication status of a secondary study from being implicitly transferred to all evidence synthesized within it.

\subsection{Classification Protocol}

For each eligible study, the classification protocol records the study type, primary object of evaluation, quality dimensions, evaluation methods, evaluation protocol, experimental regime, and level of integration.

Quality dimensions are coded independently. A component is marked as evaluated only when its quality is explicitly measured or analyzed. The mere presence of a component in the system architecture is insufficient. Thus, using retrieval does not by itself constitute retrieval-quality evaluation, using prompts does not imply input-quality evaluation, and using an LLM judge does not imply evaluation of judge reliability.

Evaluation methods and their targets are also coded separately. In particular, the use of LLM-as-a-Judge is classified as an evaluation method. Evaluator/judge quality is coded only when the evaluator itself is investigated, for example through consistency, bias, calibration, robustness, agreement, or alignment analyses.

\subsection{Levels of Evaluation Integration}

We distinguish four levels of integration.

\begin{description}
\item[None:] The study evaluates an isolated quality dimension or component without explicitly relating it to other stages.
\item[Cross-stage:] Two or more lifecycle stages or quality dimensions are explicitly related, compared, or analyzed. Examples include examining the relationship between prompt variation and output quality or between retrieval quality and answer quality.
\item[Black-box end-to-end:] A multi-stage system or workflow is evaluated from input to final outcome without decomposing the contribution of internal stages.
\item[Lifecycle-integrated:] At least three lifecycle dimensions or stages are jointly evaluated, including either operational/system quality or evaluator reliability, and the analysis preserves explicit relationships or attribution between stages.
\end{description}

This distinction prevents the phrase "end-to-end" from being treated as evidence of integrated quality assessment without considering what is actually measured.

\subsection{Human Calibration and Validation}

Because several classification decisions require semantic interpretation, the coding protocol is independently evaluated by human reviewers using stratified calibration samples. Reviewers classify studies from titles and abstracts without access to the automated classification. Cases for which the abstract is insufficient are explicitly flagged for full-text review rather than forcing a decision.

The calibration design uses partially overlapping batches so that a subset of papers is independently classified by multiple reviewers while the remaining papers increase corpus coverage. Reviewer agreement is assessed separately for categorical variables such as relevance and integration level and for individual multi-label quality dimensions. Disagreements are adjudicated and used to refine the coding protocol before large-scale application. Automated classifiers support the process but are not treated as substitutes for independent human validation.

% TODO: After the final human round, report exact sample sizes, overlap,
% agreement coefficients, disagreement handling, and full-text adjudication rate.

\section{Taxonomy of LLM-System Quality Assessment}
\label{sec:taxonomy}

The survey adopts a two-layer taxonomy. The first layer represents \emph{what is evaluated}: the lifecycle dimensions of an LLM-based system. The second represents \emph{how quality is evaluated}: the methods and protocols used to produce evidence about those dimensions. Evaluator reliability is positioned outside the normal execution lifecycle as a meta-evaluation layer.

\subsection{Lifecycle Dimensions}

The six lifecycle dimensions are summarized below.

\begin{description}

\item[Input and Specification Quality:]
Evaluates whether the task is adequately specified and whether system behavior is robust to acceptable variations in prompts, instructions, formatting, constraints, or other input-level choices. Typical concerns include ambiguity, prompt sensitivity, instruction following, and specification robustness.

\item[Context and Retrieval Quality:]
Evaluates the external information supplied to the model, including retrieval, ranking, grounding, context selection, memory retrieval, and evidence coverage. Typical concerns include relevance, recall, precision, context sufficiency, grounding, and retrieval efficiency.

\item[Output Quality:]
Evaluates the generated response or artifact. Typical concerns include correctness, factuality, faithfulness, relevance, completeness, coherence, helpfulness, safety, code correctness, and task-specific performance.

\item[Conversation and Interaction Quality:]
Evaluates behavior across multiple turns or interactive sessions. Typical concerns include consistency, state maintenance, adaptation to user feedback, dialogue-policy adherence, long-horizon behavior, and user-goal completion.

\item[Agent, Process, and Trajectory Quality:]
Evaluates intermediate decisions and execution behavior in agentic systems. Typical concerns include planning, decomposition, routing, tool selection, tool correctness, intermediate reasoning, trajectory quality, coordination, and progress toward the task goal.

\item[Operational and System Quality:]
Evaluates whether the system can operate acceptably under deployment constraints. Typical concerns include latency, throughput, monetary cost, token consumption, memory and compute usage, energy consumption, availability, privacy, safety, and observability.

\end{description}

\subsection{Cross-Cutting Properties}

Robustness, safety and security, fairness and bias, multimodality, domain specificity, and cost/efficiency are treated as cross-cutting properties rather than lifecycle stages. A robustness study, for example, may concern prompts, retrieval, outputs, agents, or evaluators. Recording the lifecycle target and the cross-cutting property separately preserves this distinction.

\subsection{Evaluation-Method Families}

Evaluation methods are classified independently of the dimension they assess:

\begin{enumerate}
    \item deterministic and rule-based evaluation;
    \item lexical or reference-based metrics;
    \item semantic or reference-aware metrics;
    \item learned evaluators or classifiers;
    \item human evaluation;
    \item LLM-as-a-Judge;
    \item multi-judge or ensemble evaluation;
    \item task- or environment-based verification;
    \item adversarial or red-team evaluation;
    \item perturbation or mutation testing;
    \item trace-aware evaluation;
    \item statistical reliability or meta-evaluation;
    \item hybrid approaches.
\end{enumerate}

The same method can be used at different lifecycle stages. For example, an LLM judge may score output factuality, assess a plan, evaluate a tool-use trajectory, or compare alternative responses. Accordingly, method frequency is not interpreted as evidence that a particular quality dimension is well covered.

\subsection{Evaluator Reliability as a Meta-Evaluation Layer}

Evaluator reliability is treated separately because it describes the quality of the measurement instrument. Relevant properties include inter- and intra-evaluator consistency, agreement with humans, calibration, bias, prompt sensitivity, presentation effects, robustness to adversarial or persuasive inputs, and dependence among evaluators.

This distinction is particularly important for LLM-as-a-Judge. The statement that a study \emph{uses} an LLM judge does not imply that it evaluates the reliability of that judge. The latter requires explicit evidence about the judge itself.

\subsection{Conceptual Overview}

Figure~\ref{fig:taxonomy} will present the taxonomy as a circular lifecycle. The inner structure represents the six system-quality dimensions. Evaluation methods form a second layer that can act on any dimension. The outer layer represents evaluator reliability and other meta-evaluation concerns.

\begin{figure*}[t]
    \centering
    \fbox{
        \parbox{0.92\textwidth}{
        \centering
        \textbf{Placeholder for final circular taxonomy}\\[5pt]
        Lifecycle dimensions: Input/Specification $\rightarrow$ Context/Retrieval
        $\rightarrow$ Output $\rightarrow$ Conversation/Interaction
        $\rightarrow$ Agent/Process/Trajectory $\rightarrow$ Operational/System\\[5pt]
        Evaluation layer: rules, reference metrics, semantic metrics, learned evaluators,
        humans, LLM judges, multi-judge schemes, task/environment checks, adversarial
        evaluation, perturbation, trace-aware evaluation\\[5pt]
        Meta-evaluation layer: reliability, agreement, consistency, calibration,
        bias, robustness
        }
    }
    \caption{Conceptual organization of LLM-system quality assessment. The lifecycle defines what is evaluated; evaluation methods define how it is measured; evaluator reliability characterizes the measurement process itself.}
    \label{fig:taxonomy}
\end{figure*}

\begin{table*}[t]
\caption{Lifecycle dimensions and representative evaluation concerns.}
\label{tab:quality-dimensions}
\begin{tabular}{p{0.18\textwidth} p{0.29\textwidth} p{0.24\textwidth} p{0.21\textwidth}}
\toprule
\textbf{Dimension} & \textbf{Central question} & \textbf{Representative properties} & \textbf{Typical evidence} \\
\midrule
Input / Specification & Is the intended task adequately and robustly specified? & sensitivity, ambiguity, instruction quality, robustness & perturbation, task performance, human assessment \\
Context / Retrieval & Was useful and appropriate evidence supplied? & relevance, recall, grounding, coverage, context sufficiency & IR metrics, semantic evaluation, human/LLM assessment \\
Output / Response & Is the resulting artifact correct and useful? & correctness, factuality, faithfulness, relevance, safety & reference metrics, humans, judges, execution tests \\
Conversation / Interaction & Does quality persist across turns and feedback? & consistency, state, adaptation, goal completion & simulation, humans, task success, dialogue traces \\
Agent / Process / Trajectory & Was the task executed through an acceptable process? & planning, tools, intermediate decisions, coordination & traces, environment checks, judges, task outcomes \\
Operational / System & Can the system operate acceptably in practice? & latency, throughput, cost, availability, resources, privacy & telemetry, profiling, monitoring, load tests \\
Meta-evaluation & Can the quality judgment itself be trusted? & agreement, calibration, bias, consistency, robustness & human comparison, reliability statistics, perturbation \\
\bottomrule
\end{tabular}
\end{table*}


% TODO: evidence-driven synthesis sections after the corpus audit.

\bibliographystyle{ACM-Reference-Format}
\bibliography{references}

\end{document}
